% ============================================================
%  Chapitre 5 — Évaluation de la fiabilité du système
% ============================================================
\chapter{Évaluation de la fiabilité du système}
\label{ch:evaluation}

\section{Introduction}

L'évaluation constitue la cinquième phase de CRISP-DM et, dans un projet RAG,
la plus délicate. Un modèle de classification se juge à son \eng{accuracy} :
on compare une classe prédite à une classe attendue. Un système génératif produit
du texte libre, pour lequel il n'existe pas de comparaison exacte à une vérité de
terrain --- deux formulations différentes peuvent être toutes deux correctes.

Ce chapitre présente donc une évaluation à deux niveaux. Nous exposons d'abord
les mesures \textbf{quantitatives de la composante de recherche}, obtenues
expérimentalement, qui ont mis en évidence une limitation majeure du système et
conduit à une préconisation étayée. Nous présentons ensuite le
\textbf{protocole d'évaluation RAGAS} de la chaîne complète, ainsi que les
métriques retenues et l'état d'avancement de cette mesure.

\section{Évaluation de la recherche sémantique}

\subsection{Protocole}

La qualité de la recherche a été mesurée indépendamment de la génération. Cette
séparation est méthodologiquement importante : si la recherche échoue à remonter
le bon document, aucune qualité de génération ne peut compenser cette
défaillance. Le contexte transmis au modèle sera erroné, et la réponse le sera
également --- fût-elle parfaitement fidèle au contexte reçu.

Sept questions ont été soumises au système, réparties en deux catégories : quatre
portant sur des informations effectivement présentes dans le corpus, et trois
portant sur des informations absentes (procédures RH, notes de frais,
intégration des nouveaux employés). Pour chacune, nous avons relevé le score de
pertinence maximal, le score minimal parmi les quatre extraits retournés, ainsi
que la latence de recherche.

\subsection{Résultats}

Les mesures sont consignées au tableau~\ref{tab:retrieval}.

\begin{table}[htbp]
\centering
\caption{Scores de pertinence mesurés sur l'index de 2\,929 segments}
\label{tab:retrieval}
\small
\begin{tabular}{@{}p{6.6cm}ccc@{}}
\toprule
\textbf{Question} & \textbf{Score max} & \textbf{Score min} & \textbf{Latence} \\
\midrule
\multicolumn{4}{@{}l}{\itshape Informations présentes dans le corpus} \\
Quels produits contient la commande 10248 ?      & 0,082 & 0,077 & 193 ms \\
Quel est le montant total de la facture 10248 ?  & 0,115 & 0,099 & 160 ms \\
Quelles boissons étaient en stock en juillet 2016 ? & 0,194 & 0,163 & 202 ms \\
\texttt{purchase orders de Maria Larsson}        & 0,567 & 0,550 & 148 ms \\
\midrule
\multicolumn{4}{@{}l}{\itshape Informations absentes du corpus} \\
Quelle est la politique de congés payés ?        & 0,046 & 0,027 & 263 ms \\
Comment soumettre une note de frais ?            & 0,003 & $-0{,}022$ & 151 ms \\
Quelle est la procédure d'intégration ?          & $-0{,}048$ & $-0{,}068$ & 158 ms \\
\bottomrule
\end{tabular}
\end{table}

\begin{figure}[htbp]
\centering
\begin{tikzpicture}[x=10cm,y=0.62cm]
  \draw[->,muted] (-0.1,0) -- (0.62,0) node[below=4pt,font=\scriptsize] {Score de pertinence maximal};
  \foreach \x/\label in {-0.1/{-0,10},0/{0},0.2/{0,20},0.4/{0,40},0.6/{0,60}} {
    \draw[muted] (\x,0.08) -- (\x,-0.08);
    \node[below=9pt,font=\scriptsize] at (\x,0) {\label};
  }
  \draw[accent,fill=accent!75] (0,0.45) rectangle (0.082,0.80);
  \draw[accent,fill=accent!75] (0,1.45) rectangle (0.115,1.80);
  \draw[accent,fill=accent!75] (0,2.45) rectangle (0.194,2.80);
  \draw[accent,fill=accent!75] (0,3.45) rectangle (0.567,3.80);
  \draw[warn,fill=warn!75] (0,4.45) rectangle (0.046,4.80);
  \draw[warn,fill=warn!75] (0,5.45) rectangle (0.003,5.80);
  \draw[warn,fill=warn!75] (-0.048,6.45) rectangle (0,6.80);
  \foreach \y/\label in {
    0.625/{Commande 10248},
    1.625/{Facture 10248},
    2.625/{Stock juillet},
    3.625/{Maria Larsson},
    4.625/{Congés payés},
    5.625/{Note de frais},
    6.625/{Intégration}
  } {
    \node[anchor=east,font=\scriptsize] at (-0.05,\y) {\label};
  }
  \node[anchor=west,font=\scriptsize] at (0.59,0.625) {0,082};
  \node[anchor=west,font=\scriptsize] at (0.59,1.625) {0,115};
  \node[anchor=west,font=\scriptsize] at (0.59,2.625) {0,194};
  \node[anchor=west,font=\scriptsize] at (0.59,3.625) {0,567};
  \node[anchor=west,font=\scriptsize] at (0.59,4.625) {0,046};
  \node[anchor=west,font=\scriptsize] at (0.59,5.625) {0,003};
  \node[anchor=west,font=\scriptsize] at (0.59,6.625) {-0,048};
\end{tikzpicture}
\caption{Scores maximaux de pertinence selon le type de question}
\label{fig:scores-retrieval}
\end{figure}

Deux constats se dégagent immédiatement.

Le premier est favorable : la \textbf{latence de recherche est très faible},
182~ms en moyenne sur les sept requêtes. L'index HNSW de ChromaDB remplit
parfaitement son office ; la recherche n'est en rien le goulot d'étranglement du
système.

Le second est préoccupant : les scores des questions \emph{présentes} dans le
corpus sont anormalement bas. Trois des quatre questions légitimes obtiennent un
score maximal inférieur à 0,20, à peine supérieur aux scores des questions
\emph{hors corpus} qui, elles, oscillent autour de 0. Un écart aussi faible
entre les deux catégories interdit tout seuillage fiable : aucune valeur de seuil
ne permettrait de rejeter les questions hors corpus sans rejeter simultanément la
majorité des questions légitimes.

\subsection{Diagnostic : l'asymétrie linguistique}

L'examen des documents effectivement retournés révèle l'origine du problème. La
question « Quels produits contient la commande 10248 ? » remonte les fichiers
\texttt{invoice\_10729.pdf} et \texttt{invoice\_10840.pdf} --- et non
\texttt{invoice\_10248.pdf}. La seule requête à obtenir un score élevé (0,567)
est aussi la seule formulée majoritairement en anglais.

L'hypothèse est donc que le modèle \texttt{all-MiniLM-L6-v2}, entraîné
essentiellement sur de l'anglais, projette mal les questions françaises dans
l'espace sémantique des documents anglais. Pour la vérifier, une expérience
contrôlée a été conduite : cinq questions ont été formulées deux fois, une fois
en français et une fois en anglais, avec un contenu sémantique identique, puis
soumises au \emph{même} index. Les résultats sont consignés au
tableau~\ref{tab:langue}.

\begin{table}[htbp]
\centering
\caption{Expérience contrôlée --- effet de la langue de la question sur le score
de pertinence (index et corpus identiques)}
\label{tab:langue}
\small
\begin{tabular}{@{}p{5.4cm}ccp{3.6cm}@{}}
\toprule
\textbf{Contenu de la question} & \textbf{Score FR} & \textbf{Score EN} &
\textbf{Gain} \\
\midrule
Produits de la commande 10248        & 0,082 & 0,337 & $+0{,}255$ \\
Montant total de la facture 10248    & 0,115 & 0,424 & $+0{,}309$ \\
Boissons en stock en juillet 2016    & 0,194 & 0,629 & $+0{,}435$ \\
Bons de commande de Maria Larsson    & 0,205 & 0,555 & $+0{,}350$ \\
Produits commandés par Paul Henriot  & 0,169 & 0,359 & $+0{,}190$ \\
\midrule
\textbf{Moyenne}                     & \textbf{0,153} & \textbf{0,461} &
\textbf{$+0{,}308$ ($\times 3{,}0$)} \\
\bottomrule
\end{tabular}
\end{table}

\begin{figure}[htbp]
\centering
\begin{tikzpicture}[x=8cm,y=0.62cm]
  \draw[->,muted] (0,0) -- (0.72,0) node[below=4pt,font=\scriptsize] {Score moyen};
  \foreach \x/\label in {0/{0},0.2/{0,20},0.4/{0,40},0.6/{0,60}} {
    \draw[muted] (\x,0.08) -- (\x,-0.08);
    \node[below=9pt,font=\scriptsize] at (\x,0) {\label};
  }
  \draw[accent,fill=accent!75] (0,0.45) rectangle (0.153,0.80);
  \draw[accent,fill=accent!75] (0,1.45) rectangle (0.461,1.80);
  \node[anchor=east,font=\scriptsize] at (-0.05,0.625) {Français};
  \node[anchor=east,font=\scriptsize] at (-0.05,1.625) {Anglais};
  \node[anchor=west,font=\scriptsize] at (0.475,0.625) {0,153};
  \node[anchor=west,font=\scriptsize] at (0.475,1.625) {0,461};
\end{tikzpicture}
\caption{Comparaison des scores moyens selon la langue de la question}
\label{fig:langue-scores}
\end{figure}

Le résultat est sans ambiguïté : à contenu sémantique constant, la simple
traduction de la question en anglais \textbf{triple le score de pertinence}
(0,461 contre 0,153 en moyenne). L'effet est également qualitatif : la question
sur les boissons en stock, formulée en anglais, remonte correctement un rapport
de stock (\texttt{StockReport\_2017-07\_1.pdf}) alors que sa version française
remontait un bon de commande, document d'un type entièrement inadapté.

\begin{encadre}
\textbf{Conclusion du diagnostic.} La cause de la faiblesse des scores n'est ni
le corpus, ni le découpage, ni le paramètre $k$, mais le \textbf{choix du modèle
de vectorisation}. Un modèle monolingue anglais est structurellement inadapté à
un usage où les questions sont posées en français sur des documents anglais.
La préconisation qui en découle est de remplacer
\texttt{all-MiniLM-L6-v2} par \texttt{paraphrase-multilingual-MiniLM-L12-v2},
entraîné sur plus de cinquante langues et conçu précisément pour aligner les
formulations translinguistiques dans un espace commun. Ce remplacement suppose
une ré-indexation complète du corpus, le changement de modèle invalidant
l'intégralité des vecteurs existants.
\end{encadre}

Ce résultat illustre l'intérêt de la démarche d'évaluation quantitative : la
faiblesse identifiée n'était pas décelable à l'usage. Les réponses produites
paraissaient acceptables ; seule la mesure des scores a révélé que le système
s'appuyait fréquemment sur les mauvais documents.

\subsection{Limite structurelle : la recherche sur identifiant exact}

Une seconde limite, indépendante de la langue, est apparue lors de ces mesures.
Même en anglais, une question portant sur un identifiant précis
(« \texttt{order 10248} ») ne remonte pas systématiquement le document portant
cet identifiant.

Ce comportement est inhérent à la recherche vectorielle. Les
\eng{embeddings} encodent le \emph{sens} d'un texte, non ses chaînes de
caractères exactes. Or les identifiants numériques \texttt{10248},
\texttt{10249} et \texttt{10729} sont sémantiquement quasi indiscernables : ce
sont tous des numéros de commande dans un contexte identique. Leur proximité
vectorielle est donc très élevée, et le classement entre eux devient
essentiellement arbitraire.

La solution reconnue pour ce type de besoin est la \textbf{recherche hybride},
combinant une composante lexicale exacte (BM25) et la composante vectorielle,
puis fusionnant les deux classements. La composante lexicale trouve
\texttt{10248} de manière exacte ; la composante vectorielle apporte la
compréhension du sens. Cette amélioration figure parmi les perspectives
retenues en conclusion.

\section{Évaluation de la chaîne complète : le protocole RAGAS}

\subsection{Les quatre métriques retenues}

L'évaluation de la chaîne complète repose sur le \eng{framework}
\textbf{RAGAS} (\eng{Retrieval-Augmented Generation Assessment}), qui décompose
la qualité d'un système RAG en dimensions mesurables séparément. Son principe
est d'employer un modèle de langue comme \emph{juge} : il compare la réponse
produite, le contexte récupéré et la réponse de référence. Quatre métriques ont
été retenues.

\paragraph{\eng{Faithfulness} (fidélité).}
C'est la métrique centrale du projet : elle mesure directement le taux
d'hallucination. La réponse est décomposée en affirmations élémentaires, et l'on
vérifie combien d'entre elles sont effectivement soutenues par le contexte
récupéré :
\[
  \text{Faithfulness} =
  \frac{\lvert\text{affirmations soutenues par le contexte}\rvert}
       {\lvert\text{affirmations totales de la réponse}\rvert}
\]
Un score de 1 signifie que rien n'a été inventé. Un score de 0,7 signifie que
30~\% des affirmations produites ne figurent pas dans les documents --- c'est-à-dire
sont des hallucinations.

\paragraph{\eng{Answer relevancy} (pertinence de la réponse).}
Elle mesure si la réponse adresse effectivement la question posée. Une réponse
peut être parfaitement fidèle au contexte tout en étant hors sujet. La métrique
procède par génération inverse : à partir de la réponse, le juge produit des
questions plausibles, puis mesure leur similarité cosinus moyenne avec la
question d'origine :
\[
  \text{Answer relevancy} =
  \frac{1}{n}\sum_{i=1}^{n}\cos\!\left(E(q),\, E(q_i')\right)
\]
où $q$ est la question initiale et $q_i'$ les questions régénérées.

\paragraph{\eng{Context precision} (précision du contexte).}
Elle évalue la composante de recherche : parmi les segments récupérés, quelle
proportion est réellement pertinente ? Elle pénalise le bruit --- les extraits
hors sujet qui diluent le contexte et augmentent le risque que le modèle
s'appuie sur un passage inadapté.

\paragraph{\eng{Context recall} (rappel du contexte).}
Elle évalue la complétude de la recherche : l'information nécessaire pour
répondre a-t-elle été intégralement retrouvée ? C'est la seule des quatre
métriques qui exige une réponse de référence rédigée par un humain, puisqu'il
faut connaître l'information attendue pour juger si elle a été retrouvée.

\vspace{0.3cm}

Ces quatre métriques se répartissent naturellement en deux couples, comme
l'indique le tableau~\ref{tab:metriques} : deux mesurent la
\textbf{génération}, deux mesurent la \textbf{recherche}. Cette séparation est
précieuse en pratique, car elle permet d'attribuer une défaillance au bon
composant.

\begin{table}[htbp]
\centering
\caption{Répartition des métriques RAGAS par composant évalué}
\label{tab:metriques}
\small
\begin{tabularx}{\textwidth}{@{}p{3.7cm}p{2.5cm}X@{}}
\toprule
\textbf{Métrique} & \textbf{Composant} & \textbf{Ce qu'un score faible indique} \\
\midrule
\eng{Faithfulness} & Génération &
Le modèle invente : revoir l'invite système ou abaisser la température \\
\eng{Answer relevancy} & Génération &
Le modèle répond à côté : reformuler l'invite \\
\eng{Context precision} & Recherche &
Trop de bruit récupéré : réduire $k$, ajouter un ré-ordonnancement \\
\eng{Context recall} & Recherche &
Information manquante : augmenter $k$, revoir le découpage ou le modèle
d'\eng{embeddings} \\
\bottomrule
\end{tabularx}
\end{table}

\subsection{Mise en œuvre}

Conformément à la contrainte de confidentialité, le modèle juge est lui aussi
exécuté localement : RAGAS est configuré pour utiliser le même Mistral 7B servi
par Ollama, et le même modèle d'\eng{embeddings}, plutôt que l'API OpenAI
qu'il emploie par défaut.

\begin{lstlisting}[style=python,caption={Configuration de RAGAS avec un juge local (\texttt{src/evaluation/ragas\_eval.py})}]
judge_llm = ChatOllama(
    base_url=settings.ollama_base_url,
    model=settings.ollama_model,
    temperature=0,
)
judge_embeddings = HuggingFaceEmbeddings(model_name=settings.embedding_model)

result = evaluate(
    dataset,
    metrics=[faithfulness, answer_relevancy, context_precision, context_recall],
    llm=judge_llm,
    embeddings=judge_embeddings,
)
\end{lstlisting}

Le jeu d'évaluation doit être fourni sous forme d'un fichier CSV à deux colonnes,
\texttt{question} et \texttt{ground\_truth}, cette dernière contenant la réponse
de référence rédigée par un relecteur humain. Le module exécute alors le pipeline
complet sur chaque question, assemble le quadruplet
\{question, réponse, contexte, référence\} attendu par RAGAS, puis calcule les
quatre métriques et enregistre le rapport dans
\texttt{docs/ragas\_report.json}.

\begin{lstlisting}[style=bash,caption={Lancement de l'évaluation}]
python -m src.evaluation.ragas_eval --testset data/eval/testset.csv
\end{lstlisting}

\subsection{État d'avancement et résultats}

\begin{encadre}
\textbf{Avertissement méthodologique.} Le module d'évaluation est développé et
opérationnel, mais la campagne de mesure RAGAS \textbf{n'a pas encore été
exécutée} à la date de rédaction de ce rapport. Deux prérequis restent à lever :
l'installation des dépendances \texttt{ragas} et \texttt{datasets} dans
l'environnement, et surtout la \textbf{constitution du jeu d'évaluation
annoté} --- une vingtaine de questions accompagnées de leurs réponses de
référence, rédigées manuellement à partir des documents du corpus. Le fichier
\texttt{data/eval/testset.csv} ne contient à ce jour qu'un gabarit de trois
questions non renseignées, dont le contenu porte de surcroît sur des documents
absents du corpus (\emph{cf.} section~\ref{sec:limite-corpus}). Les valeurs du
tableau~\ref{tab:ragas} sont donc laissées à compléter ; elles ne sont pas
estimées, afin de ne pas présenter comme mesuré ce qui ne l'a pas été.
\end{encadre}

\begin{table}[htbp]
\centering
\caption{Résultats RAGAS --- à compléter après exécution de la campagne de mesure}
\label{tab:ragas}
\small
\begin{tabular}{@{}lccp{4.3cm}@{}}
\toprule
\textbf{Métrique} & \textbf{Score obtenu} & \textbf{Cible} &
\textbf{Interprétation} \\
\midrule
\eng{Faithfulness}      & \aremplir{à mesurer} & $> 0{,}90$ &
Taux d'ancrage documentaire ; métrique critique \\
\eng{Answer relevancy}  & \aremplir{à mesurer} & $> 0{,}80$ &
Adéquation à la question posée \\
\eng{Context precision} & \aremplir{à mesurer} & $> 0{,}70$ &
Qualité du classement des extraits \\
\eng{Context recall}    & \aremplir{à mesurer} & $> 0{,}70$ &
Complétude de la recherche \\
\bottomrule
\end{tabular}
\end{table}

Les valeurs cibles indiquées correspondent aux seuils usuellement retenus dans
la littérature pour un système RAG considéré comme exploitable en production. Le
seuil le plus exigeant porte sur la \eng{faithfulness} : en contexte
d'entreprise, une réponse sur dix contenant une affirmation inventée constitue
déjà un niveau de risque difficilement acceptable.

Compte tenu du diagnostic établi à la section précédente, une prévision peut
néanmoins être formulée : les métriques de \emph{recherche}
(\eng{context precision} et \eng{context recall}) seront vraisemblablement les
plus faibles, puisque nous avons établi expérimentalement que la recherche
remonte fréquemment des documents inadaptés lorsque la question est posée en
français. La \eng{faithfulness} pourrait en revanche rester élevée : le modèle
peut fort bien rester parfaitement fidèle à un contexte... qui n'est pas le bon.
Ce cas de figure illustre précisément pourquoi une métrique unique serait
trompeuse et pourquoi la décomposition en quatre dimensions est nécessaire.

\section{Performance et latence}

Le tableau~\ref{tab:perf} présente les mesures de performance relevées.

\begin{table}[htbp]
\centering
\caption{Mesures de performance du système}
\label{tab:perf}
\small
\begin{tabularx}{\textwidth}{@{}p{6.1cm}rX@{}}
\toprule
\textbf{Étape} & \textbf{Durée} & \textbf{Remarque} \\
\midrule
Recherche vectorielle (top-$k$ = 4) & 182 ms &
Moyenne sur 7 requêtes ; négligeable \\
Chargement du modèle d'\eng{embeddings} & environ 2 s &
Une seule fois, mis en cache \\
Génération de la réponse (CPU) & environ 110 s &
Mesure relevée en conditions réelles \\
Ingestion complète du corpus & plusieurs minutes &
2\,929 segments, opération hors ligne \\
\bottomrule
\end{tabularx}
\end{table}

\begin{figure}[htbp]
\centering
\begin{tikzpicture}[x=0.022cm,y=0.72cm]
  \draw[->,muted] (0,0) -- (290,0) node[below=4pt,font=\scriptsize] {Latence (ms)};
  \foreach \x in {0,50,100,150,200,250} {
    \draw[muted] (\x,0.08) -- (\x,-0.08);
    \node[below=9pt,font=\scriptsize] at (\x,0) {\x};
  }
  \draw[accent,fill=accent!75] (0,0.45) rectangle (193,0.80);
  \draw[accent,fill=accent!75] (0,1.45) rectangle (160,1.80);
  \draw[accent,fill=accent!75] (0,2.45) rectangle (202,2.80);
  \draw[accent,fill=accent!75] (0,3.45) rectangle (148,3.80);
  \draw[warn,fill=warn!75] (0,4.45) rectangle (263,4.80);
  \draw[accent,fill=accent!75] (0,5.45) rectangle (151,5.80);
  \draw[accent,fill=accent!75] (0,6.45) rectangle (158,6.80);
  \foreach \y/\label in {
    0.625/{Commande 10248},
    1.625/{Facture 10248},
    2.625/{Stock juillet},
    3.625/{Maria Larsson},
    4.625/{Congés payés},
    5.625/{Note de frais},
    6.625/{Intégration}
  } {
    \node[anchor=east,font=\scriptsize] at (-150,\y) {\label};
  }
  \node[anchor=west,font=\scriptsize] at (268,0.625) {193};
  \node[anchor=west,font=\scriptsize] at (235,1.625) {160};
  \node[anchor=west,font=\scriptsize] at (244,2.625) {202};
  \node[anchor=west,font=\scriptsize] at (190,3.625) {148};
  \node[anchor=west,font=\scriptsize] at (268,4.625) {263};
  \node[anchor=west,font=\scriptsize] at (193,5.625) {151};
  \node[anchor=west,font=\scriptsize] at (200,6.625) {158};
\end{tikzpicture}
\caption{Latence de recherche pour les sept requêtes évaluées}
\label{fig:latence-recherche}
\end{figure}

Le déséquilibre est frappant : la recherche représente moins de 0,2~\% du temps
de réponse total, la génération plus de 99~\%. L'inférence d'un modèle de
7 milliards de paramètres sur CPU constitue l'unique goulot d'étranglement du
système.

Cette latence d'environ deux minutes est acceptable en démonstration mais
excessive pour un usage quotidien. Trois leviers permettraient de la réduire :
l'exécution sur GPU, qui apporterait typiquement un facteur dix ; le recours à un
modèle plus fortement quantifié (\texttt{mistral:7b-instruct-q4}) ou plus petit
(\texttt{phi-3}, \texttt{llama3.2:3b}) ; ou la diffusion progressive de la
réponse (\eng{streaming}), qui n'améliore pas le temps total mais réduit
considérablement la latence perçue en affichant les premiers mots dès leur
génération.

\section{Leviers anti-hallucination : bilan}

Le tableau~\ref{tab:antihallu} récapitule les mécanismes mis en place contre
l'hallucination et leur état.

\begin{table}[htbp]
\centering
\caption{Leviers de réduction des hallucinations}
\label{tab:antihallu}
\small
\begin{tabularx}{\textwidth}{@{}p{4.5cm}p{2.4cm}X@{}}
\toprule
\textbf{Levier} & \textbf{État} & \textbf{Commentaire} \\
\midrule
Invite système restrictive & Implémenté &
Interdit explicitement de sortir du contexte et fournit une formulation
d'ignorance prête à l'emploi \\
Température basse (0,1) & Implémenté &
Réduit la variabilité de génération et rend les réponses reproductibles \\
Garde-fou sur contexte vide & Implémenté &
Court-circuite le modèle si la recherche ne remonte rien \\
Citation systématique des sources & Implémenté &
Permet à l'utilisateur de vérifier chaque réponse ; levier le plus efficace en
pratique \\
Affichage des scores de pertinence & Implémenté &
Signale visuellement à l'utilisateur un appui documentaire faible \\
Seuil de score minimal & \aremplir{à implémenter} &
Écarté à ce stade : le tableau~\ref{tab:retrieval} montre qu'aucun seuil n'est
exploitable avant le passage à un modèle multilingue \\
Ré-ordonnancement (\eng{re-ranking}) & \aremplir{à implémenter} &
Un \eng{cross-encoder} améliorerait la \eng{context precision} \\
Mesure RAGAS & \aremplir{à exécuter} &
Le module est prêt ; le jeu annoté reste à constituer \\
\bottomrule
\end{tabularx}
\end{table}

Il convient de relever que le levier le plus efficace n'est pas d'ordre
algorithmique. C'est la \textbf{citation systématique des sources} : en exposant
les extraits utilisés et leur score, le système transforme une réponse
invérifiable en une affirmation contrôlable. Même si le modèle hallucine,
l'utilisateur dispose des moyens de s'en apercevoir. Cette transparence par
construction est, en contexte professionnel, plus précieuse qu'un gain marginal
de fidélité.

\section{Conclusion}

Ce chapitre a présenté l'évaluation du système à deux niveaux. L'évaluation
quantitative de la recherche a établi un résultat majeur : la latence de
recherche est négligeable (182~ms), mais les scores de pertinence sont
anormalement faibles sur des questions françaises. Une expérience contrôlée a
identifié la cause --- l'asymétrie linguistique entre un modèle
d'\eng{embeddings} anglophone et des questions françaises --- et l'a quantifiée :
la traduction de la question en anglais triple le score de pertinence. Cette
mesure fonde la préconisation d'adopter un modèle multilingue. Une seconde limite
structurelle a été documentée : l'incapacité de la recherche vectorielle à
retrouver un identifiant exact, qui appelle une recherche hybride.

Le protocole d'évaluation RAGAS de la chaîne complète est développé et
opérationnel, mais la campagne de mesure reste à conduire : elle requiert la
constitution préalable d'un jeu de questions annotées. Les emplacements
correspondants ont été laissés explicitement vacants plutôt que renseignés par
estimation. Le chapitre suivant présente le déploiement du système.
